\thesisabstract{%
Questa tesi di laurea magistrale studia modelli di intelligenza multimodale consapevoli dell'affidabilità per sistemi cyber-fisici. La demo combina una struttura orientata alla ricerca, domande esplicite, una metodologia riproducibile, una valutazione quantitativa e una discussione della validità. Il sistema visivo utilizza la palette ISISLab, accenti oro UniSA controllati, trattamenti trasparenti del logo e componenti editoriali pensati per la scrittura tecnica estesa. Sostituire il testo con una sintesi strutturata di motivazione, lacuna scientifica, metodo, risultati principali, contributo originale e implicazioni.
}

\declarationpage{%
Dichiaro che il presente elaborato è frutto del mio lavoro originale e che tutte le fonti, i dati, il software e i contributi esterni sono stati correttamente riconosciuti. Adattare questa dichiarazione ai requisiti ufficiali dell'Università degli Studi di Salerno e al regolamento del corso di laurea magistrale.
}

\acknowledgementspage{%
Utilizzare questa pagina facoltativa per ringraziare la supervisione scientifica, i colleghi di laboratorio, i collaboratori, gli enti finanziatori, i fornitori dei dataset e le persone che hanno offerto supporto personale. Rimuovere la pagina quando non è richiesta.
}
